% ---------------------------------------------------------------------
%  第 11 章 物理观测量
% ---------------------------------------------------------------------
\chapter{物理观测量}
\label{ch:observables}

\cref{tab:observables}按物理量列出观测接口及所需输入。

Green 函数本身不是可观测量，但它包含全部单粒子信息。本章说明库
如何从中提取密度矩阵、动能、关联能、电流与谱函数。所有函数都只读
输入，不修改容器；它们都接受 0-based 的时间指标。

\section{密度矩阵}
\label{sec:density}

密度矩阵由小于分量的等时值给出：
\begin{equation}
    \rho = \ii\,\xi\,\Gless(t,t),
    \qquad
    \xi = -1\ (\text{费米型}),\ +1\ (\text{玻色型}).
\end{equation}
对费米型，等时小于分量为 $\Gless(t,t)=\ii\langle c^{\dagger}c\rangle$，
代入得 $\rho=\langle c^{\dagger}c\rangle$，即占据数。
在平衡态，它也可以由 Matsubara 分量给出：
\begin{equation}
    \rho = -\Gmat(\bet).
\end{equation}
\code{density_matrix} 同时支持这两种入口：\code{tstp == MATSUBARA}
时取 Matsubara 块，否则取 $t$ 处的等时块。\code{density} 返回密度矩阵
的迹，即总占据数 $N(t)$。

\paragraph{其他容器。}
同一观测量在三个变体上都有对应函数：
\begin{itemize}
    \item \code{pseudo_density} / \code{pseudo_density_matrix}：赝粒子
          对象，前者返回迹（占据约束 $n \to 1$ 的监控量），后者返回一个块；
    \item \code{pseudo_density_matrix_step}：赝粒子时间步视图；
    \item \code{steady_density} / \code{steady_density_matrix}：取稳态
          对象的零相对时间块，前者返回迹。
\end{itemize}
赝粒子接口也使用小于分量的等时块，统计因子仍为 $\ii\xi$。

\section{动能}
\label{sec:kinetic}

动能由哈密顿量与密度矩阵给出：
\begin{equation}
    E_{\mathrm{kin}}(t) = \mathrm{Re}\,\Tr\big[\bm{H}(t)\,\rho(t)\big].
\end{equation}
常数与含时情形分别使用固定哈密顿量块和当前时间步的块。

\paragraph{截断版本。}
接口 \code{kinetic_energy_trunc} 使用两个移动窗对象 $G$ 与 $D$，
并对当前领头行的历史块进行积分。以 $X_l$ 表示对象 $X$ 的相对索引 $(0,l)$ 块，
\begin{equation}
    E_{\mathrm{kin}}(t)
    =
    h\,\mathrm{Im}\,\Tr\Big[\sum^{t_c}_{l=0}w_l
    \big((G_l^R)^{\dagger}D_l^<-(G_l^<)^{\dagger}D_l^R\big)\Big],
\end{equation}
其中 $w_l$ 是两端 Gregory 权重、中间为 $1$ 的求积权重。该函数
\code{kinetic_energy_trunc} 要求 $t_c\ge 2k+2$，因为两端的修正
带不能重叠。

\section{关联能}
\label{sec:corr-energy}

关联能由自能与 Green 函数的卷积给出：
\begin{equation}
    E_{\mathrm{corr}}(t)
    =
    \tfrac12\,\mathrm{Re}\,\Tr\big[\rho_{\Sig G}(t)\big],
    \qquad
    \rho_{\Sig G} = \texttt{convolution\_density\_matrix\_herm}(\Sig,G).
\end{equation}
实现复用第\ref{ch:convolution}章的等时卷积：先得到
$\rho_{\Sig G}$，再求迹并取实部。因子 $1/2$ 来自自能的对称化。

\section{电流}
\label{sec:current}

对多轨道或多 $k$ 点体系，电流是各通道的加权和：
\begin{equation}
    J = \sum_{k} w_k\,
    \mathrm{Re}\,\Tr\big[\bm{v}_k\,\rho_k\big],
\end{equation}
其中 $\bm{v}_k$ 是速度算符，$\rho_k$ 是第 $k$ 个通道的密度矩阵，
$w_k$ 是权重（例如 $k$ 点的求积权重）。函数 \code{current}
接收三个等长数组，逐项求迹后加权求和。

\section{Wigner 切片}
\label{sec:wigner}

\cref{fig:wigner}标出了固定平均时间时提取切片的方向与范围。

Wigner 切片把双时间 Green 函数投影到“平均时间 + 相对时间”坐标。
给定平均时间索引 $n_{\mathrm{av}}$，令 $t_{\mathrm{av}}=n_{\mathrm{av}}h$。
索引偏移 $j$ 对应的物理时间为
\begin{equation}
    t_1 = t_{\mathrm{av}}+jh,\qquad t_2 = t_{\mathrm{av}}-jh,\qquad
    t_{\mathrm{rel}} = 2jh,
\end{equation}
其中 $j=0,\dots,\min(n_{\mathrm{av}},\,n_t-n_{\mathrm{av}})$。函数
\code{wigner_slices} 返回三样东西：相对时间数组、推迟分量数组与
小于分量数组，后两者的每一项是一个 $d\times d$ 块。

\begin{figure}[htbp]
    \centering
    \begin{tikzpicture}[scale=0.85, >=Stealth]
        % axes
        \draw[->, thick] (0,0) -- (7.2,0) node[right, font=\small] {$t'$};
        \draw[->, thick] (0,0) -- (0,7.2) node[above, font=\small] {$t$};
        % diagonal
        \draw[dashed, gray] (0,0) -- (6.6,6.6);
        \node[font=\scriptsize, gray, anchor=south east] at (6.0,6.0)
            {$t=t'$};
        % slice through t_av = 3 (perpendicular to the diagonal)
        \draw[very thick, dashed, gray] (5.3,0.7) -- (3,3);
        \draw[very thick, kborange, ->] (3,3) -- (0.7,5.3);
        \fill[kborange] (0.7,5.3) circle (2.4pt);
        \node[font=\scriptsize, kborange, anchor=south west] at (0.75,5.35)
            {$(t_{av}-s,\,t_{av}+s)$};
        \node[font=\scriptsize, gray, anchor=north east] at (5.25,0.65)
            {$t_{\mathrm{rel}}<0$};
        \fill[kbblue] (3,3) circle (2.2pt);
        \node[font=\scriptsize, kbblue, anchor=north west] at (3.05,2.95)
            {$s=0$};
        % projection of t_av
        \draw[dotted, gray] (3,0) -- (3,3);
        \node[font=\scriptsize, below] at (3.1,0) {$t_{av}$};
    \end{tikzpicture}
    \caption[Wigner 切片]{Wigner 切片。固定平均时间 $t_{av}$ 后，切片垂直于对角线；
             橙色半线对应接口返回的 $t_{\mathrm{rel}}\ge0$，其中
             $s=t_{\mathrm{rel}}/2$。灰色虚线表示负相对时间方向。}
    \label{fig:wigner}
\end{figure}

\section{加窗谱函数}
\label{sec:spectrum}

谱函数由推迟分量的虚部给出：
\begin{equation}
    A(t_{av},\omega)
    =
    -\frac{1}{\pi}\,\mathrm{Im}
    \int\!\dd t_{\mathrm{rel}}\;
    \ee^{\ii\omega t_{\mathrm{rel}}}\,
    w(t_{\mathrm{rel}})\,\Gret(t_{av}+t_{\mathrm{rel}}/2,\;
    t_{av}-t_{\mathrm{rel}}/2),
\end{equation}
其中 $w$ 是窗函数，用于抑制截断振荡。函数
\code{evaluate_spectrum_windowed} 支持两种积分方法：
\begin{itemize}
    \item \code{Cubic}：三次修正的 DFT，要求区间数为偶数且大于
          $16$，否则报错而不回退；
    \item \code{Linear}：对分段线性插值做精确积分，至少需要两个点。
\end{itemize}
两者都要求相对时间网格等距，并对每个轨道对角元分别计算谱函数。

\begin{table}[htbp]
    \centering
    \caption[观测量一览]{观测量一览}
    \label{tab:observables}
    \small
    \begin{tabularx}{\textwidth}{@{}llX@{}}
        \toprule
        量 & 函数 & 输入 \\
        \midrule
        密度（迹） & \code{density} &
            轮廓对象与时间步，或 Matsubara 标记 \\
        密度矩阵 & \code{density_matrix} &
            轮廓对象与时间步，或 Matsubara 标记 \\
        密度（迹） & \code{steady_density} & 稳态对象 \\
        密度矩阵 & \code{steady_density_matrix} &
            稳态对象 \\
        密度（迹） & \code{pseudo_density} & 赝粒子对象 \\
        密度矩阵 & \code{pseudo_density_matrix} &
            赝粒子对象 \\
        动能 & \code{kinetic_energy} / \code{kinetic_energy_cf} &
            哈密顿量与轮廓对象 \\
        动能 & \code{kinetic_energy_trunc} &
            移动窗对象 \\
        关联能 & \code{correlation_energy} &
            自能与 Green 函数 \\
        电流 & \code{current} &
            速度算符、密度矩阵与权重的数组 \\
        Wigner 切片 & \code{wigner_slices} &
            轮廓对象与平均时间 \\
        谱函数 & \code{evaluate_spectrum_windowed} &
            推迟切片、时间网格、频率与窗函数 \\
        \bottomrule
    \end{tabularx}
\end{table}

\section{单位与因子}
\label{sec:units}

所有观测量都在库的内部单位下计算：$\hbar=1$、$k_B=1$。时间以
步长 $h$ 的整数倍表示，能量以 $1/h$ 的量纲出现。转换到
物理单位需要两个因子：
\begin{itemize}
    \item 时间：每个时间步对应 $h$（由调用者决定）；
    \item 能量：$\bm{H}$ 与 $\mu$ 的数值直接以能量为单位，
          因此由 $\bm{H}$ 构造的动能自然带有正确的量纲。
\end{itemize}
密度矩阵是无量纲的占据数；电流的量纲由速度算符决定。报告不引入
额外的单位换算层，以免在链条中引入隐藏因子。
